강한 연결 요소(strongly connected component, SCC)는 방향 그래프에서 정점 $u$에서 정점 $v$로 가는 경로와 정점 $v$에서 정점 $u$로 가는 경로가 모두 존재하는 정점들의 집합이다. 즉, 강한 연결 요소는 방향 그래프에서 서로 도달할 수 있는 정점들의 집합이다. 강한 연결 요소를 이용하면 임의의 방향 그래프를 방향 비순환 그래프로 압축할 수 있다. 이번 장에서는 방향 그래프에서 강한 연결 요소를 찾는 알고리즘을 다룬다.

\section{강한 연결 요소}
그래프에 사이클이 있다면 유향 그래프여도 사이클에 포함된 모든 정점 사이를 도달하는 경로가 존재한다. 따라서 이러한 정점들의 집합을 하나로 묶어보자는 아이디어를 낼 수 있다. 수학적으로 엄밀히 서술하면 다음과 같다:

\begin{definition}[강한 연결 요소]
    유향 그래프 $G=(V, E)$의 강한 연결 요소는 $G$의 부분 그래프이며, 이 부분 그래프의 모든 정점 쌍 $(u, v)$에 대해 $u$에서 $v$로 가는 경로와 $v$에서 $u$로 가는 경로가 모두 존재하는 그래프를 말한다.
\end{definition}

강한 연결 요소를 구하는 가장 쉬운 방법은 $O(V^2)$의 시간복잡도를 가지는 알고리즘을 사용하는 것이다. 모든 정점에 대해 DFS 또는 BFS를 수행하여 도달할 수 있는 정점을 찾고, 그 정점에서 다시 DFS 또는 BFS를 수행하여 원래의 정점에 도달할 수 있는지 확인하자. 도달할 수 있다면 SCC에 포함시키면 되고, 아니라면 새로운 SCC를 만들면 된다. 이 방법은 모든 정점에 대해 DFS 또는 BFS를 수행하므로 시간복잡도는 $O(V(V+E))$가 된다.

하지만 더 효율적인 알고리즘이 존재한다. 대표적으로 Kosaraju 알고리즘과 Tarjan 알고리즘이 있다. 두 알고리즘 모두 DFS를 기반으로 하여 $O(V+E)$의 시간복잡도로 강한 연결 요소를 찾을 수 있다.

\section{Kosaraju 알고리즘}
코사라주 알고리즘은 두 번의 DFS를 이용하여 강한 연결 요소를 찾는 알고리즘이다. 첫 번째 DFS에서는 각 정점의 종료 시간을 기록하고, 두 번째 DFS에서는 그래프의 모든 간선의 방향을 뒤집은 후 종료 시간이 가장 늦은 정점부터 DFS를 수행하여 강한 연결 요소를 찾는다. 현재까지 방문하지 않은 정점들 중에서, 뒤집어진 그래프에서 방문 가능한 정점을 모으면 하나의 강한 연결 요소가 된다.

이 방법으로 SCC를 구할 수 있는 이유를 간단하게 생각해보자. 먼저,
\begin{enumerate}
    \item 어떤 정점 $v$에서 DFS를 수행한다고 하자. 그러면, $v$와 같은 SCC에 속한 정점과, $v$에서 도달할 수 있는 다른 SCC에 속한 정점들을 모두 방문하게 된다.
    \item 이제 거꾸로 뒤집은 그래프의 $v$에서 DFS를 수행한다고 하자.
    \begin{itemize}
        \item 원래 그래프의 $v$에서 도달할 수 없었던 점들은 더 늦은 종료 시간을 갖는다.
        \item 뒤집어진 그래프의 $v$에서 도달할 수 있는 점들은 원래 그래프의 $v$에서 도달할 수 없었던 점들이거나 같은 SCC에 속한 점들이다.
        \item 따라서 뒤집어진 그래프의 $v$에서 DFS를 수행하면, $v$와 같은 SCC에 속한 정점들만 모두 방문하게 된다. 왜냐하면 원래 그래프에서 $v$에서 도달할 수 없었던 점들은 더 늦은 종료 시간을 가져 이미 두 번째 과정에서 DFS를 수행하면서 방문했기 때문이다.
    \end{itemize}
\end{enumerate}

이 내용을 다음과 같이 엄밀히 수학적 귀납으로 증명할 수 있다.
\begin{proof}
    (보충)
\end{proof}

구현 시에 종료 시간을 기록하여도 되지만, 먼저 종료된 정점부터 스택에 넣어도 충분하다. 구현 코드는 \Cref{code: Kosaraju}을 참고하라.

\begin{code} \label{code: Kosaraju}
Kosaraju 알고리즘으로 SCC를 찾는 코드
\begin{lstlisting}
#include <bits/stdc++.h>
using namespace std;
int n, m;
vector<int> adj1[202020], adj2[202020];

int ch[202020], val[101010];
stack<int> stk;
vector< vector<int> > scc;
void dfs1(int v) {
    ch[v] = 1;
    for (int i = 0; i < adj1[v].size(); i++) {
        int u = adj1[v][i];
        if (ch[u]) continue;
        dfs1(u);
    }
    stk.push(v);
}

void dfs2(int v) {
    ch[v] = 1;
    scc.back().push_back(v);
    for (int i = 0; i < adj2[v].size(); i++) {
        int u = adj2[v][i];
        if (ch[u]) continue;
        dfs2(u);
    }
}

int main() {
    // make adj1 and adj2
    for (int i = 1; i <= n; i++) ch[i] = 0;
    for (int i = 1; i <= n; i++) {
        if (ch[i]) continue;
        dfs1(i);
    }
    for (int i = 1; i <= n; i++) ch[i] = 0;
    while (stk.size() > 0) {
        int v = stk.top();
        stk.pop();
        if (ch[v]) continue;
        scc.push_back(vector<int>());
        dfs2(v);
    }
    // scc contains all strongly connected components
}
\end{lstlisting}
\end{code}

시간복잡도는 DFS가 지배하므로 $O(V+E)$이다.

\section{Tarjan 알고리즘}
(설명, 증명, 구현 추가)

\section{2-SAT 문제}
2-SAT 문제는 강한 연결 요소를 이용하여 해결할 수 있는 문제이다. 2-SAT 문제는 다음과 같이 정의된다:

\section{연습문제 A}
% \begin{problem}
%     최소 스패닝 트리로 가능한 것 하나를 출력하는 프로그램을 작성하라.
% \end{problem}
% \begin{problem}
%     최소 스패닝 트리로 가능한 것이 하나인지, 여러 개인지 판별하는 프로그램을 작성하라.
% \end{problem}
% \begin{problem}
%     $N$개의 수로 이루어진 수열 $A_i$가 주어진다. 여러분은 $i$번 정점과 $j$번 정점을 비용 $\gcd(A_{i+1}, A_{i+2}, \cdots, A_j)$로 연결할 수 있다. $1$번부터 $N$번까지의 정점이 포함된 그래프 $G$를 연결 그래프로 만들기 위해 필요한 최소 비용을 $O(N^2)$의 시간복잡도로 구하는 프로그램을 작성하라. 단, $\gcd$는 최대공약수 함수이다.
% \end{problem}

\section{연습문제 B}

\subsection{기초문제}
% \begin{problem} \url{https://www.acmicpc.net/problem/1197} \end{problem}
% \begin{problem} \url{https://www.acmicpc.net/problem/4343} \end{problem}
% \begin{problem} \url{https://www.acmicpc.net/problem/4386} \end{problem}
% \begin{problem} \url{https://www.acmicpc.net/problem/9134} \end{problem}
% \begin{problem} \url{https://www.acmicpc.net/problem/18769} \end{problem}
% \begin{problem} \url{https://www.acmicpc.net/problem/20010} \end{problem}
% \begin{problem} \url{https://www.acmicpc.net/problem/24010} \end{problem}
% \begin{problem} \url{https://www.acmicpc.net/problem/25545} \end{problem}
% \begin{problem} \url{https://www.acmicpc.net/problem/27009} \end{problem}
% \begin{problem} \url{https://www.acmicpc.net/problem/28019} \end{problem}

\subsection{응용문제}
% \begin{problem} \url{https://www.acmicpc.net/problem/1396} \end{problem}
% \begin{problem} \url{https://www.acmicpc.net/problem/3900} \end{problem}
% \begin{problem} \url{https://www.acmicpc.net/problem/7788} \end{problem}
% \begin{problem} \url{https://www.acmicpc.net/problem/10335} \end{problem}
% \begin{problem} \url{https://www.acmicpc.net/problem/12667} \end{problem}
% \begin{problem} \url{https://www.acmicpc.net/problem/14574} \end{problem}
% \begin{problem} \url{https://www.acmicpc.net/problem/18919} \end{problem}
% \begin{problem} \url{https://www.acmicpc.net/problem/20263} \end{problem}
% \begin{problem} \url{https://www.acmicpc.net/problem/23108} \end{problem}
% \begin{problem} \url{https://www.acmicpc.net/problem/23736} \end{problem}

\subsection{도전문제}
% \begin{problem} \url{https://www.acmicpc.net/problem/4203} \end{problem}
% \begin{problem} \url{https://www.acmicpc.net/problem/5257} \end{problem}
% \begin{problem} \url{https://www.acmicpc.net/problem/8527} \end{problem}
% \begin{problem} \url{https://www.acmicpc.net/problem/11592} \end{problem}
% \begin{problem} \url{https://www.acmicpc.net/problem/12668} \end{problem}
% \begin{problem} \url{https://www.acmicpc.net/problem/14828} \end{problem}
% \begin{problem} \url{https://www.acmicpc.net/problem/16074} \end{problem}
% \begin{problem} \url{https://www.acmicpc.net/problem/17582} \end{problem}
% \begin{problem} \url{https://www.acmicpc.net/problem/18482} \end{problem}
% \begin{problem} \url{https://www.acmicpc.net/problem/18963} \end{problem}
% \begin{problem} \url{https://www.acmicpc.net/problem/19251} \end{problem}
% \begin{problem} \url{https://www.acmicpc.net/problem/19852} \end{problem}
% \begin{problem} \url{https://www.acmicpc.net/problem/20086} \end{problem}
% \begin{problem} \url{https://www.acmicpc.net/problem/20503} \end{problem}
% \begin{problem} \url{https://www.acmicpc.net/problem/24953} \end{problem}
% \begin{problem} \url{https://www.acmicpc.net/problem/24961} \end{problem}
% \begin{problem} \url{https://www.acmicpc.net/problem/25121} \end{problem}
% \begin{problem} \url{https://www.acmicpc.net/problem/27071} \end{problem}
% \begin{problem} \url{https://www.acmicpc.net/problem/27877} \end{problem}
% \begin{problem} \url{https://www.acmicpc.net/problem/28002} \end{problem}